\chapter{AGENTS.md 设计详解} \label{app:agents-md}

本书正文第~\ref{ch:deadline-ai}~章介绍了 \texttt{AGENTS.md} 的基本概念——它相当于发给 AI 的“员工手册”。本章以一份真实项目中使用的 \texttt{AGENTS.md} 为例，逐节解释每条规则的设计意图和实际作用。需要说明的是，这份 \texttt{AGENTS.md} 是在单人开发的项目背景下定制的（团队规模仅作者一人），因此没有涉及多人协作中的确认、审批、分工等流程，而是聚焦于“如何让 AI 安全、高效地辅助一个人完成项目”。

如果您的项目规模不同（例如团队开发），可以在本章基础上增加协作相关规则；但核心思想——用明确规则替代口头约定、用自动化流程减少手误、用止损机制控制风险——是通用的。

\section{为什么单人项目更需要一份 \texttt{AGENTS.md}}

在团队中，编码规范通常靠代码审查和同事提醒来维持。但单人项目没有第二双眼睛盯着，所有风险都靠自己扛。此时，AI 的行为一致性、可回滚性、以及“不乱动”的安全边界，反而比团队项目更重要。\texttt{AGENTS.md} 就是把这些边界写下来，让 AI 在每次会话中都按同一套规则工作，避免“上次说的这次忘了”的情况。

\section{通用原则的设计意图}

\subsection{编译输出目录隔离}

\begin{lstlisting}
	- 所有编译输出必须放入 ./Codex-Debug 目录。
	- 严禁使用 Qt 默认编译目录 ./Debug。
	- 严禁在项目根目录编译。
\end{lstlisting}

单人开发时，您可能同时用 Qt Creator 手动编译和让 AI 辅助编译。如果 AI 的编译产物和您的混在一起，清理时容易误删，备份时也会把大体积的中间文件包进去。单独给 AI 划一个 \texttt{./Codex-Debug}，既方便您区分“AI 改过的版本”和“您自己改的版本”，也让备份机制只备份源码而非编译产物。

\subsection{增量备份机制}

这是单人项目最重要的安全网。AI 改代码可能很快，但改出问题也很快。增量备份只记录“每次会话首次修改前”的源码状态，不备份编译输出和临时文件，空间开销小，却能随时回滚。

设计上的几个关键点：

\begin{itemize}
	\item 排除 \texttt{./Backup}、\texttt{./Codex-Debug}、\texttt{./Result}，避免备份嵌套和冗余；
	\item 按修改时间（mtime）决定是否覆盖，只保留较新版本，不重复存储未变化的文件；
	\item 源工程已删除的文件仍然保留在备份中，回滚时可以把误删的文件找回来；
	\item \texttt{./Backup} 是增量副本，不是独立快照，这意味着它不能精确还原到某次会话开始时的状态，但足以恢复所有曾经存在过的文件的最新版本。
\end{itemize}

对于单人开发者，这套机制比 Git 更轻量——您不需要每次修改前手动提交，AI 会在首次修改前自动做一次增量更新。您只需要记住：改崩了就用 \texttt{RESTORE ALL} 或 \texttt{RESTORE PARTIAL} 回滚。

\subsection{清理临时文件}

AI 在调试过程中经常生成临时脚本或日志，然后“忘记”删除。要求 AI 每次任务完成后自己清理，如果清理失败必须报告，可以防止项目根目录被垃圾文件堆满，也避免您事后手动翻找。

\section{代码与界面规范的设计意图}

这一节是针对 Qt/C++ 项目的特定规范。如果您的项目不是 Qt，可以整体删掉或替换为对应的框架规范。这里以原文为例说明设计思路。

\subsection{国际化（i18n）}

在 Qt 中，UI 文本必须用 \texttt{tr()} 包裹才能被翻译工具提取。单人项目往往没有专职翻译，但如果您以后要发布多语言版本，现在写好 \texttt{tr()} 比事后补要省力得多。规则中“新增未同步条目记录在 \texttt{未更新语言条目.md}”是一种过渡策略——允许您先把功能写完，再集中处理翻译，避免中断编码思路。

\subsection{UI 构建：优先使用 \texttt{.ui} 文件}

单人项目尤其需要“透明”的界面描述。\texttt{.ui} 文件是 XML，您可以直接看坐标、尺寸、属性，比动态创建代码更容易理解和修改。您不需要依赖可视化设计器也能手动编辑，而且 AI 对 XML 的处理比解析 C++ 布局代码更可靠。

\subsection{WBS 开发进度维护}

在单人项目中，WBS（工作分解结构）仍然有价值，但可以大幅简化。原文中的“负责人确认”“同事反馈”“导出提醒”等内容在单人场景下可以精简为：

\begin{itemize}
	\item 维护一份简单的任务清单（如 \texttt{00\_WBS开发进度.md}），记录每个子任务的状态（未开始/进行中/已完成）；
	\item 状态更新由您自己决定，AI 只能建议，不能擅自标记完成；
	\item 不必每周导出 CSV，除非您需要自己做进度回顾。
\end{itemize}

之所以保留“AI 不能擅自标记完成”这条，是因为单人项目没有别人帮您验收，AI 可能误判“编译通过”为“任务完成”，而实际功能并未验证。所以规则中明确写道：“仅在对应范围已完成且有验证依据，或负责人（即您本人）明确确认时，才可标记完成”——在单人场景下，负责人就是您自己，确认方式可以是您口头告诉 AI“这个任务我验收过了”。

\section{测试与验证的设计意图}

\subsection{测试类型划分}

在单人项目中，您的时间更宝贵，因此用 Python 做快速验证、用 C++ 做最终确认的策略尤其有效。原文规定“Python 验证通过后再移植到 C++”虽然原文中有“若用户要求”，但单人项目中您完全可以自己决定是否移植，AI 只是按照规则执行。

\subsection{测试结果文件命名规则}

固定命名规则让您每次跑完测试后都能快速找到历史记录，方便对比不同版本的结果。单人项目没有测试团队帮您存档，全靠自己，所以这个规则很实用。

\section{止损策略的设计意图}

止损策略（STOP-LOSS）是单人项目最重要的防线。AI 可能在您没有实时监控的情况下连续多轮改错方向，而您可能直到很久之后才发现。止损策略要求：如果 AI 连续 3 轮 GREEN 测试（即逻辑失败，而非编译或运行时崩溃）都未通过，就必须停止并报告，避免无谓消耗。

边界情况例外的设计也很重要——有些理论上的 bug 实际上永远不会触发，而您可能急着交付，这时候允许 AI 在提供证据后请示您是否可以放过。举证责任在 AI 身上，您只需做最后判断。

\section{命令参考的设计意图}

把常用操作封装成命令（如 \texttt{BACKUP}、\texttt{TEST PYTHON}、\texttt{RESTORE ALL}），是为了减少您和 AI 之间的沟通成本。您只需说一句“备份”或“回滚全部”，AI 就按既定流程执行，不必每次都重复解释要做什么。这对单人开发者尤其友好——您不需要记住所有细节，只需记住几个命令名。

\section{如何根据您的项目调整这份 \texttt{AGENTS.md}}

这份 \texttt{AGENTS.md} 是针对 Qt/C++ 桌面项目写的。如果您的项目是嵌入式裸机、Python 后端或 Web 前端，可以做如下裁剪：

\begin{enumerate}
	\item \textbf{嵌入式裸机}：删除国际化、UI 构建章节，测试类型只保留 C 测试或硬件自检，编译输出目录改为 \texttt{./build} 或 \texttt{./Codex-Build}；
	\item \textbf{Python 后端}：删除 Qt UI 规范，增加 PyTest 测试规范，备份规则保留；
	\item \textbf{Web 前端}：删除 C++ 测试，增加 ESLint/Prettier 自动格式化规则。
\end{enumerate}

但无论如何裁剪，有三条原则是单人项目必须坚守的：

\begin{enumerate}
	\item 把 AI 的编译输出和您的操作目录分开；
	\item 任何 AI 对代码的修改都必须可回滚；
	\item AI 不能替您做“完成”或“验收”的最终判断。
\end{enumerate}

这三条原则，正是正文第~\ref{ch:deadline-ai}~章中“AI 是体力劳工，不是军师”在单人项目中的具体体现。您可以根据自己的项目环境，把这三条原则翻译成具体的目录名和规则条款。